% sections/03-design-rationale.tex
% OWNER: Member B
% Do NOT edit this file if you are not Member B.
% Statements marked "our interpretation" are the authors' analysis, not claims made by Uber.

\section{Design Rationale}
\label{sec:design-rationale}

This section explains \emph{why} Michelangelo looks the way it does. It lists the
problems Uber faced when it started the platform (\Cref{sec:rationale-problems}),
traces each design decision back to those problems
(\Cref{sec:rationale-decisions}) and names the quality attributes the design
targets (\Cref{sec:rationale-qualities}). \emph{How} the resulting components are
arranged into layers is covered in \Cref{sec:layered-architecture}.

% ------------------------------------------------------------
\subsection{Uber's Problems in 2015}
\label{sec:rationale-problems}

Uber ``began building Michelangelo in mid 2015''~\cite{delbalso2017meet}. At that time ML
was ``not widely used at Uber'', and the company had ``no centralized ML efforts
and a few bespoke ML systems''~\cite{delbalso2018scaling}. The 2017 post names the concrete
gaps, which we number P1--P6 for use in the rest of this section:

\begin{enumerate}[label=\textbf{P\arabic*}]
  \item \label{prob:tools}\textbf{Fragmented tools.} Data scientists ``were
        using a wide variety of tools to create predictive models''~\cite{delbalso2017meet}.
  \item \label{prob:pipelines}\textbf{No reproducible pipelines.} There were
        ``no systems in place to build reliable, uniform, and reproducible
        pipelines''~\cite{delbalso2017meet}.
  \item \label{prob:desktop}\textbf{Desktop-sized training.} It was ``not
        possible to train models larger than what would fit on data scientists'
        desktop machines''~\cite{delbalso2017meet}.
  \item \label{prob:experiments}\textbf{No experiment tracking.} There was no
        ``standard place to store the results of training
        experiments''~\cite{delbalso2017meet}.
  \item \label{prob:deploy}\textbf{No standard deployment path.} There was ``no
        established path to deploying a model into production''. For each project,
        the engineering team ``had to create a custom serving container'', and
        separate teams built ``bespoke one-off systems'' to use models in
        production~\cite{delbalso2017meet}.
  \item \label{prob:features}\textbf{Duplicated feature work.} ``Many modeling
        problems at Uber use identical or similar features''~\cite{delbalso2017meet}.
        In addition, models deployed online ``cannot access data stored in
        HDFS''~\cite{delbalso2017meet}, so online models need another way to get their
        features at serving time.
\end{enumerate}

Uber's stated goal was to close these gaps ``by standardizing the workflows and
tools across teams'' and to ``create a system that would grow with the
business''~\cite{delbalso2017meet}.

% ------------------------------------------------------------
\subsection{From Requirements to Design Decisions}
\label{sec:rationale-decisions}

\Cref{tab:rationale} turns each problem into a requirement, the decision Uber made
to meet it and the component that implements the decision. The requirement
wording and the problem mapping in the last column are \emph{our
interpretation}; the decisions and components are taken from Uber's posts.

\begin{table}[ht]
  \centering
  \small
  \caption{Requirements, design decisions and components of Michelangelo.}
  \label{tab:rationale}
  \begin{tabularx}{\linewidth}{@{}L{2.6cm}>{\raggedright\arraybackslash}X L{2.8cm}L{1.3cm}@{}}
    \toprule
    \textbf{Requirement} & \textbf{Design decision} & \textbf{Component} & \textbf{Solves} \\
    \midrule
    Reuse features across teams
      & One shared, curated feature store; offline features in Hive tables on HDFS, online
        features precomputed into Cassandra~\cite{delbalso2017meet}
      & Feature store (later Palette)~\cite{delbalso2017meet,uber2024palette}
      & \ref{prob:features}, \ref{prob:tools} \\
    \addlinespace
    Same features in training and serving
      & A feature DSL whose ``same expressions are applied at training time and
        at prediction time''~\cite{delbalso2017meet}
      & Feature DSL, shared by training and serving~\cite{delbalso2017meet}
      & \ref{prob:pipelines}, \ref{prob:features} \\
    \addlinespace
    Train beyond one machine
      & Distributed training on Spark~\cite{delbalso2017meet}, later Horovod for
        deep learning~\cite{delbalso2018scaling} and Ray
        trainers~\cite{wang2024predictive}
      & Training jobs on Spark, Horovod,
        Ray~\cite{delbalso2017meet,delbalso2018scaling,wang2024predictive}
      & \ref{prob:desktop} \\
    \addlinespace
    Store and compare every experiment
      & Store each trained model as a versioned object with its configuration,
        learned parameters and evaluation report; compare models in the
        UI~\cite{delbalso2017meet}
      & Model repository~\cite{delbalso2017meet}, later Gallery~\cite{sun2020gallery}
      & \ref{prob:experiments}, \ref{prob:pipelines} \\
    \addlinespace
    One-click, safe deployment
      & Deploy through the UI or API; at serving time a model is found by
        ``its UUID and an optional tag (or alias)''~\cite{delbalso2017meet}
      & Model repository + deployment tags~\cite{delbalso2017meet}
      & \ref{prob:deploy} \\
    \addlinespace
    No custom serving per project
      & One shared prediction service that other services call over RPC,
        instead of a container per project~\cite{delbalso2017meet}
      & Online Prediction Service~\cite{delbalso2017meet,wang2024predictive}
      & \ref{prob:deploy}, \ref{prob:tools} \\
    \addlinespace
    One workflow for all teams
      & A single end-to-end platform with a web UI, API and workflow system over
        shared infrastructure~\cite{delbalso2017meet}
      & Michelangelo platform (API tier, UI)~\cite{delbalso2017meet}
      & \ref{prob:tools}, \ref{prob:pipelines} \\
    \bottomrule
  \end{tabularx}
\end{table}

\paragraph{Why these decisions.}
\begin{itemize}
  \item \textbf{A central feature store instead of per-team features.} Because
        many models use ``identical or similar features''~\cite{delbalso2017meet},
        computing a feature once and sharing it removes duplicate work
        (\ref{prob:features}). Uber later called the feature store ``critical to
        enabling teams to reuse important predictive features''~\cite{delbalso2018scaling}.
        The split into an offline and an online store exists because online
        models cannot read HDFS~\cite{delbalso2017meet}; it later became the offline and
        online data planes (\Cref{sec:layered-mapping}).
  \item \textbf{A shared DSL instead of separate training and serving code.}
        Running the same expression in both paths is meant ``to help guarantee
        that the same final set of features is generated and sent to the model in
        both cases''~\cite{delbalso2017meet}. This targets the non-uniform pipelines of
        \ref{prob:pipelines}.
  \item \textbf{Distributed training instead of desktop tools.} The training
        system ``scales up to handle billions of samples and down to small
        datasets''~\cite{delbalso2017meet}, which removes the desktop limit of
        \ref{prob:desktop}. Why the engine could later change from Spark to Ray
        is discussed in \Cref{sec:layered-evaluation}.
  \item \textbf{A versioned model repository instead of ad hoc result files.}
        Every trained model is stored with its full configuration (features
        used, hyper-parameter values), learned parameters and evaluation
        report~\cite{delbalso2017meet}. This gives \ref{prob:experiments} a standard
        place and ties a result to the inputs that produced it
        (\ref{prob:pipelines}).
  \item \textbf{Tags and one shared serving path instead of custom containers.}
        A client asks for a model by UUID or tag~\cite{delbalso2017meet}; \emph{our
        interpretation} is that a new model can therefore replace an old one
        without any change to client code. Uber notes that this ``allows
        safe transitions from old models to new models and side-by-side A/B
        testing of models''~\cite{delbalso2017meet}. Together with one prediction
        service, this provides the missing deployment path of \ref{prob:deploy}.
\end{itemize}

% ------------------------------------------------------------
\subsection{Quality Attributes the Design Targets}
\label{sec:rationale-qualities}

The decisions above serve five quality attributes~\cite{bass2021}.
\Cref{tab:rationale-qualities} ties each one to a mechanism of Michelangelo and,
where Uber publishes one, a number.

\begin{table}[ht]
  \centering
  \small
  \caption{Quality attributes, mechanisms and published figures.}
  \label{tab:rationale-qualities}
  \begin{tabularx}{\linewidth}{@{}L{2.6cm}*{2}{>{\raggedright\arraybackslash}X}@{}}
    \toprule
    \textbf{Attribute} & \textbf{Mechanism} & \textbf{Published figure} \\
    \midrule
    Scalability
      & Distributed training on Spark~\cite{delbalso2017meet},
        Horovod~\cite{delbalso2018scaling} and Ray~\cite{wang2024predictive}; for offline prediction, ``add
        more Spark executors and let Spark manage the
        parallelism''~\cite{delbalso2017meet}
      & Training on ``billions of samples''; top models served more than
        250,000 predictions/s in 2017; 10 million real-time predictions/s at peak in
        2024~\cite{delbalso2017meet,wang2024predictive} \\
    \addlinespace
    Performance (low latency)
      & Online features precomputed into Cassandra ``where they can be read at
        low latency''; a shared online prediction service~\cite{delbalso2017meet}
      & Typical P95 latency under 5\,ms without Cassandra features and under
        10\,ms with them~\cite{delbalso2017meet} \\
    \addlinespace
    Reusability
      & One feature store where teams ``share, discover, and use a highly
        curated set of features''~\cite{delbalso2017meet}
      & About 10,000 features in the feature store in 2017~\cite{delbalso2017meet} \\
    \addlinespace
    Reproducibility
      & Versioned model objects with full configuration, parameters and
        evaluation report; one DSL for training and serving~\cite{delbalso2017meet}
      & Every trained model is stored as a versioned object~\cite{delbalso2017meet} \\
    \addlinespace
    Deployability
      & Deployment via UI or API; models addressed by UUID or tag, enabling safe
        transitions and A/B tests~\cite{delbalso2017meet}
      & No published figure; \emph{our interpretation:} a tag switch replaces
        a model with no change to calling services \\
    \bottomrule
  \end{tabularx}
\end{table}

Each attribute answers one or more 2015 problems: \emph{scalability} answers
\ref{prob:desktop} and the goal of a system that would ``grow with the
business''~\cite{delbalso2017meet}; \emph{low latency} is needed because other Uber
services call the prediction service during live requests~\cite{delbalso2017meet}, and
its per-hop cost is examined in \Cref{sec:layered-evaluation};
\emph{reusability} answers \ref{prob:features}; \emph{reproducibility} answers
\ref{prob:pipelines} and \ref{prob:experiments}, since each model version keeps
the features and hyper-parameters that produced it~\cite{delbalso2017meet}; and
\emph{deployability} answers \ref{prob:deploy}, in line with Uber's lesson to
``[m]ake the deployment process fast and magical''~\cite{delbalso2018scaling}.
